\subsection{Sampling}

We introduce a regex-parameterized sampling function $\Sample$ in Algorithm~\ref{alg:sampling}. $\Sample$ gives a recursive means of converting any regex $R$ into an efficient sampling procedure, whose random outputs are (for unambiguous $R$) unbiased samples from the conditional u-MPS distribution associated with the subset $R \subset \Sigma^*$. This is formalized in
%
\begin{theorem}
\label{thm:sample}
Consider a u-MPS model with core tensor $\mc{A}$ and boundary vectors $\alpha$ and $\omega$, along with an unambiguous regex $R$ whose right transfer operator $\ER_R$ converges. Let $P$ indicate the probability distribution over arbitrary strings defined by the u-MPS, so that $\Sigma_{s\in\Sigma^*} P(s) = 1$. Then calling $\Sample(R, \alphamat, \omegamat)$ samples a string $s \in \Sigma^*$ from the conditional u-MPS distribution $P(s | s \in R) = P(s) / P(R)$, with $s \in R$ and where $P(R) := \sum_{s' \in R} P(s')$.
\end{theorem}
%
We prove Theorem~\ref{thm:sample} in Appendix~\ref{sec:appendix_b}, which also discusses the use of ambiguous regex $R$. For this latter case, Algorithm~\ref{alg:sampling} works identically, but weights strings $s$ by the number of times $s$ matches $R$.

Although Algorithm~\ref{alg:sampling} is written in a recursive manner, it is useful to consider the simple example $R = \Sigma^n$, a concatenation of the single-character regex $\Sigma$ with itself $n$ times, to understand the overall control flow. In this case, Algorithm~\ref{alg:sampling} first attempts to sample the initial character in the string via a recursive call to $\Sample(\Sigma, \alphamat, \ER_{\Sigma^{n-1}}(\omegamat))$. This requires $n-1$ applications of the transfer operator $\ER$ to the initial right boundary matrix, and yields one new character before continuing to the right and repeating this process again.

As is common with recursive algorithms, caching intermediate information permits the naive cost of $(n-1) + (n-2) + \cdots + 1 = \bigO(n^2)$ transfer operator applications to be reduced to $\bigO(n)$. This cached version is equivalent to a simple iterative algorithm, where a sequence of right boundary matrices is first generated and saved during a right-to-left sweep, before a left-to-right sweep is used to sample text and propagate conditional information using the left boundary matrices. Using this idea, we show in Appendix~\ref{sec:appendix_c} that for typical regex $R$, Algorithm~\ref{alg:sampling} can be run with average-case runtime $\bigO(L d D^3)$ and worst-case memory usage $\bigO(L D^2)$, for $L$ the number of characters in $R$, $d$ the size of $\Sigma$, and $D$ the bond dimension of the u-MPS.

